\section{应用神经网络：量化损失}

\subsection{示例问题：能否通过考试？}

Amini 用一个直觉性的例子来说明神经网络的应用流程。问题是：根据学生的出勤次数 $x_1$ 和在期末项目上花费的小时数 $x_2$，预测学生是否能通过 6.S191 课程。

\begin{figure}[H]
\centering
\includegraphics[width=0.85\textwidth]{slides-images/slide-047.jpg}
\caption{示例问题：给定新数据点 $[4,5]$，预测是否通过？\protect\footnotemark}
\end{figure}
\footnotetext{Slides 第 47 页。}

将数据点 $x^{(1)} = [4, 5]$ 送入一个具有单隐藏层的网络，假设网络（使用随机初始化的权重）输出 $\hat{y}_1 = 0.1$。但该学生实际通过了课程（真实标签 $y^{(1)} = 1$）。这说明当前的网络权重还不够好——我们需要一种方法来量化预测与真实值之间的差距。

\subsection{损失函数}

\textbf{损失函数}（Loss Function）衡量的是网络预测值与真实值之间的差距。

对于单个数据点的损失：
$$\mathcal{L}\big(f(x^{(i)}; \mathbf{W}), y^{(i)}\big)$$

\begin{figure}[H]
\centering
\includegraphics[width=0.85\textwidth]{slides-images/slide-050.jpg}
\caption{损失函数：衡量预测与真实值之间的偏差\protect\footnotemark}
\end{figure}
\footnotetext{Slides 第 50 页。}

整个数据集上的\textbf{经验损失}（Empirical Loss）是所有样本损失的平均：
$$J(\mathbf{W}) = \frac{1}{n}\sum_{i=1}^{n}\mathcal{L}\big(f(x^{(i)}; \mathbf{W}), y^{(i)}\big)$$

经验损失也被称为\textbf{目标函数}（Objective Function）、\textbf{代价函数}（Cost Function）或\textbf{经验风险}（Empirical Risk）。

\begin{importantbox}{两种核心损失函数}
\textbf{1. 交叉熵损失}（Binary Cross Entropy Loss）——用于分类问题（输出为概率）：
$$J(\mathbf{W}) = -\frac{1}{n}\sum_{i=1}^{n}\left[y^{(i)}\log f(x^{(i)};\mathbf{W}) + (1-y^{(i)})\log(1-f(x^{(i)};\mathbf{W}))\right]$$

\textbf{2. 均方误差损失}（Mean Squared Error Loss）——用于回归问题（输出为连续值）：
$$J(\mathbf{W}) = \frac{1}{n}\sum_{i=1}^{n}\left(y^{(i)} - f(x^{(i)};\mathbf{W})\right)^2$$
\end{importantbox}

\begin{figure}[H]
\centering
\begin{subfigure}[t]{0.48\textwidth}
\includegraphics[width=\textwidth]{slides-images/slide-052.jpg}
\caption{交叉熵损失（分类问题）}
\end{subfigure}
\hfill
\begin{subfigure}[t]{0.48\textwidth}
\includegraphics[width=\textwidth]{slides-images/slide-053.jpg}
\caption{均方误差损失（回归问题）}
\end{subfigure}
\caption{两种核心损失函数\protect\footnotemark}
\end{figure}
\footnotetext{Slides 第 52--53 页。}

\begin{warningbox}{选择正确的损失函数}
损失函数的选择必须与任务类型匹配：
\begin{itemize}
\item 分类任务（Pass/Fail、猫/狗）使用交叉熵损失
\item 回归任务（预测成绩分数、预测房价）使用均方误差
\end{itemize}
错误的损失函数会导致模型无法正确学习。例如，将 MSE 用于分类任务会因为损失函数对概率边界不敏感而导致训练缓慢。
\end{warningbox}

\subsection{本章小结}

损失函数是训练神经网络的核心组件，它量化了模型预测与真实标签之间的差距。经验损失是所有训练样本上损失的平均值。分类任务通常使用交叉熵损失，回归任务使用均方误差。训练的目标就是找到使经验损失最小的权重参数。

